\section{直線と bracket が非零の枝の交点}
\subsection{交点の isotropy Lie algebra}

\begin{proposition}\label{prop:line-conic-slopes}
正規な実際の退化因子の被約特異支持に直線 $L$ と滑らかな枝 $\Gamma$ が含まれ、
両者の一般横断型がともに $\widetilde E_8$ であり、異なる接線方向で $p$ において交わるとする。
$\Gamma$ の conormal bracket は $p$ で非零であると仮定する。
このとき両枝の modular residue の $p$ における一次係数は零である。
\end{proposition}

\begin{proof}
両枝を局所的に $x$ 軸と $y$ 軸に取り、normal 二方向を残す。
modular slope をそれぞれ $a,b$ とする。
補題\ref{lem:zero-slope-nilpotent}により、一方が零で他方が非零ということはない。
従って矛盾のため $a,b\ne0$ と仮定する。

原点の isotropy Lie algebra を $\mathfrak g=T_p^*\PP^4$、両枝の conormal ideals を
$H_x,H_y$ とし、$N=H_x\cap H_y$ と置く。$N$ は二次元の ideal である。
一次 Poisson tensor は二つの normal 座標だけに依存するので
$[\mathfrak g,\mathfrak g]\subset N$。
$Z(p)=0$ による unimodularity から $N$ は abelian となる。
具体的には $N=\langle u,v\rangle$ で $[u,v]=\alpha u+\beta v$ とすると、
$\operatorname{tr}\operatorname{ad}(u)=\beta$、
$\operatorname{tr}\operatorname{ad}(v)=-\alpha$ なので両者が零である。

まず直線側の conormal bracket も $p$ で非零とする。
このとき $H_x,H_y$ はともに Heisenberg algebra。
$\operatorname{ad}(dx)|_N$ と $\operatorname{ad}(dy)|_N$ は非零の二次冪零行列であり、
Jacobi と $N$ の可換性により可換である。従って両者の像と核は同じ一次元空間 $K$。
これは二つの Heisenberg ideals の共通の中心である。
leading-center weight の式 $\nu=-(k+1)a/2$ を $k=0$ に適用すると、
$K$ の modular weight は $-a/2$ と $-b/2$ の両方であり、$a=b$。
しかし tangent 二軸は独立な eigenvectors なので eigenvalue $a$ の重複度は二以上となり、
\eqref{eq:line-spectrum-node}に反する。

次に直線側の conormal bracket が $p$ で零とする。
$H_x$ は abelian、$H_y$ は Heisenberg である。
$M=\operatorname{ad}(dx)|_N$ は非零の二次冪零行列で、その像と核の共通の一次元空間を
$K_0$ と書く。$K_0$ は $H_y$ の中心なので、その modular weight は $\nu=-b/2$。

$K=\ker(\operatorname{ad}(dx):H_x\to H_x)$ と置くと
\[
 K\cap N=K_0,\qquad K/K_0\hookrightarrow H_x/N.
\]
従って $K$ に現れる modular eigenvalues は $\nu$ と、存在するなら $b$ だけである。
この kernel は modular 線形化 $D=Z_1^*$ により保存される。
実際、$D(dx)=a\,dx+n$、$n\in N$ であり、$D$ は isotropy algebra の derivation だから、
$v\in H_x$ に対し
\[
 D([dx,v])=a[dx,v]+[n,v]+[dx,Dv]
 =a[dx,v]+[dx,Dv]
\]
となる。ここでは $H_x$ の可換性を用いた。

直線 conormal pencil の零点は単純なので、その leading center の modular weight は $-a$。
また $\widetilde E_8$ の tangent Hamiltonian の中心固有値が $-R/2$ であることから、
この leading center は $K$ に含まれる。従って
\[
 -a\in\{b,-b/2\},\qquad b/a\in\{-1,2\}.
\]
(LS2) は比二を許さないため、$b=-a$。
$M$ は $N/K_0$ を $K_0$ に非零に写し、derivation $D$ に対する weight shift は $a$ なので、
$N/K_0$ の weight は $\nu-a=-a/2$、$K_0$ の weight は $\nu=a/2$ となる。
よって ambient spectrum は
\[
 a\{1,-1,-1/2,1/2\}.
\]
これは $m=0,1,2$ のいずれの (LS2) とも一致しない。矛盾である。
\end{proof}

\paragraph{conic の非消滅条件と二配置の排除。}
別途証明した conic-plane fourth-jet 補題により、normal quintic の $\widetilde E_8$ conic
では平面内 normal 方向の Hessian は一般に非零である。
conormal splitting $\OO(-4)\oplus\OO(-2)^2$ で bracket matrix の対応する $B_{00}$ は
定数であり、その非零性により conormal bracket は conic 全体で非零である。
従って上の命題を各 conic--line 交点で適用できる。

$211a$（conic が鎖の中央）では conic 上の二つの異なる交点がともに二重 modular 零点
となり、$\PP^1$ 上の非零 vector field の零因子次数二に反する。
$211b$（直線が鎖の中央）でも中央の直線は conic との交点で二重零を持ち、
さらに別の直線との交点で零を持つので、同様に矛盾する。
